\section{Introduction}
Let $M$ be a smooth connected manifold. Then $M$ is {\it locally
  homogeneous} if, for some homogeneous space $G/H$, the smooth
structure of $M$ can be stif\/fened to a ${\mathcal G}$-structure, where
${\mathcal G}$ is the pseudogroup of all those local transformations
of $G/H$ that are restrictions of a left translation by an element of~$G$.

\subsection{Main results}\label{goose}%
The chief purpose of this article is to re-examine local homogeneity
from the Lie groupoid point of view. This leads, in particular, to the
conclusion that a locally homogeneous manifold can be
inf\/initesimalized to obtain a transitive Lie algebroid over $M$,
equipped with a f\/lat linear connection $\nabla$ that is compatible
with the Lie algebroid structure, i.e.\ is a f\/lat {\it Cartan}
connection (see Section~\ref{skoj} below). As we shall elucidate, $\nabla$
being f\/lat amounts to the existence of a~transitive action by some Lie
{\em algebra} on~$M$ `twisted' by a monodromy representation.

Not all transitive Lie algebroids equipped with a f\/lat Cartan
connection are inf\/initesimali\-za\-tions of a locally homogeneous
manifold. If $({\mathfrak g}, \nabla)$ is an inf\/initesimalization,
then $\nabla $ must be {\it geometrically closed}, in a sense made
precise below. Fortunately, geometric closure is suf\/f\/icient for
reversing the inf\/initesimalization procedure:

\begin{TheoremBlaom}\label{TheoremA}
  A smooth manifold $M$ is locally homogeneous if and only if there
  exists a transitive Lie algebroid $\mathfrak g$ over $M$ admitting a
  flat, geometrically closed, Cartan connection $\nabla$.
\end{TheoremBlaom}


As we show in Section~\ref{main}, the particular model $G/H$ that applies is
encoded in the torsion and monodromy of $\nabla$. By contrast, in the
predominant approach to local homogeneity, one f\/ixes a particular
homogeneous space $G/H$ {\em a priori}, and asks if $M$ admits $G/H$
as a local model. In practice, this requires one to anticipate an
appropriate model, or check several candidates systematically. See,
e.g., the survey~\cite{Goldman_10} for this point of view.

\subsection{Geometric closure}\label{geomclose}

Let ${\mathfrak g} $ be a transitive Lie algebroid equipped with a
f\/lat Cartan connection $\nabla$. Fix an arbitrary point $m_0 \in M$
and an arbitrary simply-connected open neighbourhood $U$ of $m_0$. Let
${\mathfrak g}_0$ be the f\/inite-dimensional vector space of $\nabla
$-parallel sections over $U$. Because $\nabla$ is f\/lat, ${\mathfrak
  g}_0$ is the same for all choices of $U$, up to obvious
identif\/ications. Because $\nabla$ is Cartan, ${\mathfrak g}_0 \subset
\Gamma({\mathfrak g}_U)$ is a~subalgebra (with bracket encoded in the
torsion of $\nabla $; see Section~\ref{localform}).

Denote the simply-connected Lie group having ${\mathfrak g}_0 $ as its
Lie algebra by $G_0$, and let ${\mathfrak h}_0$ denote the kernel of
the map
\begin{gather*}
  \xi \mapsto \#\xi (m_0) \colon \  {\mathfrak g}_0 \rightarrow
  T_{m_0}M.
\end{gather*}
Here $\# \colon {\mathfrak g} \rightarrow TM$ denotes the anchor of
$\mathfrak g$. It is easy to see that ${\mathfrak h}_0 \subset
{\mathfrak g}_0$ is a subalgebra, and we say $\nabla $ is {\it
  geometrically closed} if the connected subgroup $H_0 \subset G_0$
with Lie algebra ${\mathfrak h}_0$ is closed in $G_0$. In making this
def\/inition, the choice of f\/ixed point $m_0 \in M$ is immaterial, as we
will establish later in Section~\ref{hudo}.

The basic prototype of a Lie algebroid supporting a f\/lat Cartan
connection is the action algebroid ${\mathfrak g}_0 \times M$
associated with some inf\/initesimal action of a Lie {\em algebra}
${\mathfrak g}_0$ on $M$ (see Section~\ref{morph2}). Indeed, locally this is
the only example (see Section~\ref{localform}). In this case, ${\mathfrak
  h}_0 \subset {\mathfrak g}_0$ is the isotropy subalgebra at $m_0$.

For an example a Lie algebra action that is {\em not} geometrically
closed, see \cite[Example~8]{Kamber_Michor_04} and~\cite{Kowalski_97}.


